\documentclass[12pt,a4paper]{article}

% ── Packages ──────────────────────────────────────────────────────────────────
\usepackage[a4paper, margin=2.5cm]{geometry}
\usepackage[T1]{fontenc}
\usepackage[utf8]{inputenc}
\usepackage{lmodern}
\usepackage{amsmath, amssymb}
\usepackage{booktabs}
\usepackage{array}
\usepackage{tabularx}
\usepackage{longtable}
\usepackage{parskip}
\usepackage{hyperref}
\usepackage{xcolor}
\usepackage{titlesec}
\usepackage{enumitem}
\usepackage{listings}
\usepackage{microtype}
\usepackage{fancyhdr}
\usepackage{lastpage}
\usepackage{graphicx}
\usepackage{float}

% ── Hyperref setup ────────────────────────────────────────────────────────────
\hypersetup{
    colorlinks=true,
    linkcolor=blue!70!black,
    urlcolor=blue!70!black,
    citecolor=blue!70!black,
    pdftitle={COS 710 Assignment 3 -- Structure-Based Grammatical Evolution for Regression},
    pdfauthor={Yohali Malaika Kamangu}
}

% ── Header / Footer ───────────────────────────────────────────────────────────
\pagestyle{fancy}
\fancyhf{}
\rhead{COS 710 -- Assignment 3}
\lhead{Yohali Malaika Kamangu (u23618583)}
\cfoot{\thepage\ of \pageref{LastPage}}
\renewcommand{\headrulewidth}{0.4pt}

% ── Section title formatting ──────────────────────────────────────────────────
\titleformat{\section}{\large\bfseries}{}{0em}{\thesection.\quad}
\titleformat{\subsection}{\normalsize\bfseries}{}{0em}{\thesubsection\quad}

% ── Code listing style ────────────────────────────────────────────────────────
\lstset{
    basicstyle=\ttfamily\small,
    breaklines=true,
    frame=single,
    backgroundcolor=\color{gray!8},
    rulecolor=\color{gray!40},
    xleftmargin=1em,
    xrightmargin=1em
}

% ── Column type for tables ────────────────────────────────────────────────────
\newcolumntype{L}[1]{>{\raggedright\arraybackslash}p{#1}}
\newcolumntype{C}[1]{>{\centering\arraybackslash}p{#1}}

% ── Numbered References (make thebibliography use \section, not \section*) ────
\makeatletter
\renewcommand\thebibliography[1]{%
  \section{\refname}%
  \@mkboth{\MakeUppercase\refname}{\MakeUppercase\refname}%
  \list{\@biblabel{\@arabic\c@enumiv}}%
       {\settowidth\labelwidth{\@biblabel{#1}}%
        \leftmargin\labelwidth
        \advance\leftmargin\labelsep
        \@openbib@code
        \usecounter{enumiv}%
        \let\p@enumiv\@empty
        \renewcommand\theenumiv{\@arabic\c@enumiv}}%
  \sloppy\clubpenalty4000\widowpenalty4000%
  \sfcode`\.\@m}
\makeatother

% ──────────────────────────────────────────────────────────────────────────────
\begin{document}

% ── Title Page ────────────────────────────────────────────────────────────────
\begin{titlepage}
    \centering
    \vspace*{2cm}
    {\LARGE\bfseries COS 710: Artificial Intelligence\par}
    \vspace{0.5cm}
    {\Large Assignment 3: Structure-Based Grammatical Evolution for Regression\par}
    \vspace{1.5cm}
    {\large Yohali Malaika Kamangu\par}
    {\large u23618583\par}
    \vspace{0.5cm}
    {\large May 2026\par}
    \vspace{1.5cm}
    {\normalsize Dataset: Residential Energy Usage UK (2014--2020)\par}
\end{titlepage}

\tableofcontents
\newpage

% ──────────────────────────────────────────────────────────────────────────────
\section{Description of the Data Set}

The data set is the Residential Energy Usage United Kingdom data set
covering the years 2014--2020. Each row in the raw CSV is one observation
sampled at 15--minute intervals and includes the columns
\texttt{utc\_timestamp}, \texttt{Electricity\_load},
\texttt{Residential\_electricity\_price},
\texttt{Residential\_solar\_generation},
\texttt{Residential\_wind\_generation}, \texttt{Temperature} and
\texttt{Relative Humidity}. Only the \texttt{Electricity\_load} column is
used in this assignment, in keeping with Assignments~1 and~2: the task is
defined as a univariate time--series regression problem in which the
target $y_t$ is the load value at time $t$ and the inputs are the seven
most recent observed lag values $(y_{t-1}, y_{t-2}, \ldots, y_{t-7})$.

To keep run--times tractable but the data sample large enough for
stable statistics, the first $100{,}000$ valid rows of the CSV are read.
A sliding window of size 7 is then applied to produce
$99{,}993$ supervised examples. These are split temporally
(no shuffling) into $79{,}994$ training and $19{,}999$ test
examples ($80/20$). Temporal ordering is preserved because shuffling a
time--series would leak future information into the training set.

% ──────────────────────────────────────────────────────────────────────────────
\section{Description of the Performance Metric}

Mean Squared Error (MSE) is used as the primary measure of predictor
quality:
\begin{equation}
\mathrm{MSE} \;=\; \frac{1}{n}\sum_{i=1}^{n}\bigl(\hat{y}_i - y_i\bigr)^{2}
\end{equation}
where $\hat{y}_i$ is the value produced by the evolved expression on the
$i$--th windowed input and $y_i$ the actual measured load. MSE is chosen
for three reasons: (i) it is differentiable and convex, providing a smooth
selection pressure surface; (ii) it penalises large errors quadratically,
which matches the operational cost of grid--load mis--prediction;
(iii) it is identical to the metric used in Assignments~1 and~2, ensuring
that the three sets of results are directly comparable.

Two secondary indicators are also reported: the \emph{hits count}
(training cases with $|\hat{y}_i-y_i| < 0.01$) and the \emph{success
predicate} (a run is considered successful if its training MSE is below
$2 \times 10^{-4}$, the same threshold used in A1 and A2).

% ──────────────────────────────────────────────────────────────────────────────
\section{How Structure is Incorporated into Canonical GE}

Canonical Grammatical Evolution \cite{ONeillRyan2001} represents each
individual as a variable--length integer codon genome which is mapped
through a BNF grammar to a phenotype program. Selection, crossover and
mutation operate on the genome (the genotype); fitness is computed on
the phenotype. Canonical GE has no explicit structural component: every
behaviourally good genome is selected with the same probability, even
when the population has collapsed onto a single derivation skeleton.

In this work canonical GE is extended with a \emph{structural
fitness--sharing} penalty applied at the derivation level. During mapping
the sequence of (non--terminal, production) pairs chosen by the leftmost
derivation is recorded as a \emph{derivation signature} ---
a compact representation of the program's structural skeleton that is
independent of its constant values. Two individuals are considered to
occupy the same structural niche when their normalised prefix--match
ratio exceeds a threshold $\sigma$. The raw fitness is then inflated
by a multiplicative penalty proportional to the niche count:
\begin{equation}
f_{\text{adj}}(i) \;=\; f_{\text{raw}}(i) \;\bigl(1 + \alpha \cdot
\mathrm{nicheCount}(i)\bigr)
\label{eq:sharing}
\end{equation}
with $\alpha = 0.01$ and $\sigma = 0.6$ used in this work. The niche count
is estimated from $50$ uniformly sampled peers per individual rather than
the full population to keep the per--generation cost linear in the
population size.

\paragraph{Why derivation signatures and not phenotype trees.}
The structural sharing in Assignment~2 used prefix matching on the GP
expression \emph{tree}. In GE the equivalent natural object is the
derivation tree, but the derivation tree carries strictly more
information than the phenotype tree (it records which grammar
production was chosen even when several productions lead to the same
phenotype, e.g.\ productions 0 and 1 of $\langle\text{expr}\rangle$).
This makes the GE signature a finer--grained structural measure --- two
individuals are only considered similar if the grammar walked exactly
the same path, which is the right notion of similarity for a
grammar--based search.

\paragraph{Elitism uses raw, not adjusted, fitness.}
A subtle but important design point (highlighted as feedback on
Assignment~2): the goal of elitism is to guarantee that the best
\emph{behavioural} solution discovered so far is carried unchanged into
the next generation. Whether that individual happens to share its
structural skeleton with many peers is irrelevant to its predictive
quality, so the elite slot is chosen by raw MSE--based fitness. The
adjusted fitness in Equation~\ref{eq:sharing} only influences
\emph{parent selection}, which is where structural diversity needs to be
encouraged.

% ──────────────────────────────────────────────────────────────────────────────
\section{Description of the Grammar Used}

A compact BNF grammar suitable for closed--form symbolic regression on
the seven lag features is used:
\begin{lstlisting}
<expr>  ::= <expr> <op> <expr>          [production 0]
          | ( <expr> <op> <expr> )      [production 1]
          | <var>                       [production 2]
          | <const>                     [production 3]
<op>    ::= + | - | * | /
<var>   ::= x1 | x2 | x3 | x4 | x5 | x6 | x7
<const> ::= <digit>.<digit><digit>      [single production]
<digit> ::= 0 | 1 | 2 | 3 | 4 | 5 | 6 | 7 | 8 | 9
\end{lstlisting}

\paragraph{Why these productions.} The $\langle\text{expr}\rangle$
non--terminal has four productions chosen to give the grammar both
recursive expressive power (productions 0 and 1) and a terminating route
(productions 2 and 3). Including the parenthesised production alongside
the bare one slightly biases the derivation toward sub--expressions
that are clearly delimited; functionally the two yield the same AST
shape but doubling the count of recursive options means a random codon
is twice as likely to grow the tree as to terminate it, balancing the
two terminating productions and producing a healthy depth distribution
in the initial population.

\paragraph{Why these terminals (the seven--lag input set).}
The terminal set $\{x_1,\ldots,x_7\}$ corresponds to the seven most
recent observed load values $(y_{t-1},\ldots,y_{t-7})$. Seven lags
were chosen for three concrete reasons. First, the residential energy
data is sampled at 15--minute intervals, so seven lags cover the most
recent $1\,\text{h}\,45\,\text{min}$ of load --- the sub--hourly window
over which the auto--correlation function of the series is strongest.
Second, a small autocorrelation pilot study performed for Assignment~1
on the same series showed that the seven--lag window captures over
$95\%$ of the linear predictive power and that adding further lags
produced negligible gains. Third, using the same terminal set as A1 and
A2 keeps the search space identical across the three assignments, so any
performance differences can be attributed to the search engine rather
than to the input encoding.

\paragraph{Why include ephemeral random constants.}
The $\langle\text{const}\rangle$ production yields three--digit ERCs in
$[0.00,9.99]$. Without ERCs the grammar can only construct expressions
whose numerical value comes from input variables, which makes it
impossible to represent simple but useful transformations such as
\emph{averaging} the two most recent lags (which requires the constant
$2$, or the equivalent $1.99$ produced by the best run of this
assignment). ERCs are therefore essential for the closure of the
representable function space under arithmetic scaling and offsetting.
The constant range $[0.00, 9.99]$ was chosen empirically: it covers the
small scaling factors typical of normalised time--series averaging
without permitting wildly large constants that destabilise selection.

\paragraph{Termination guarantees.}
Two safeguards prevent the mapping from looping forever. (a) The codon
stream is allowed to \emph{wrap} at most twice; a third wrap declares
the genome invalid and assigns it the worst fitness so it is removed by
selection. (b) Once the derivation depth equals six --- the same depth
cap used for GP trees in A1/A2 --- $\langle\text{expr}\rangle$ is
restricted to its terminal productions ($\langle\text{var}\rangle$ or
$\langle\text{const}\rangle$). This is the depth--bounded mapping used
in standard GE implementations such as PonyGE2.

% ──────────────────────────────────────────────────────────────────────────────
\section{Initial Population Generation}

Each genome in the initial population is constructed by drawing
$64$ uniform random integers in $[0, 256)$. This is the canonical GE
initialisation \cite{ONeillRyan2001}.

\paragraph{Why random codon initialisation and not sensible
initialisation.} GE has two well--known initialisation strategies:
\emph{random codon initialisation} (used here) and \emph{sensible
initialisation} (Azad and Ryan, 2003), which is the GE analogue of GP's
ramped half--and--half: it grows the derivation tree to controlled
depths and back--solves for the codons needed to produce it. Sensible
initialisation guarantees that no individual in generation 0 is invalid
or trivially small, which is its main advantage over random
initialisation.

Random codon initialisation was preferred in this work for three
reasons. (1)~\emph{Structural diversity for sharing.} The structural
fitness--sharing penalty (Section~3) depends on a heterogeneous initial
population: if every genome already shares a similar derivation
skeleton (as sensible initialisation tends to produce because it draws
depths from a small ramp) the sharing penalty has little to push
against in the early generations. Pure random codon initialisation
guarantees a wide spread of derivation signatures from generation~0.
(2)~\emph{Termination is already guaranteed.} The depth--bounded
fallback in the grammar (Section~4) ensures every random codon stream
maps to a complete sentence within the depth cap, removing the main
motivation for sensible initialisation. (3)~\emph{Simplicity / parity
with the literature.} Random codon initialisation is the canonical
choice in the original O'Neill and Ryan formulation and the most common
choice in published GE benchmark studies, making the results easier to
compare.

The genome length of 64 codons was chosen as a conservative upper bound
on the number of derivation steps required by an expression at the
maximum depth of 6: a balanced binary $\langle\text{expr}\rangle$ tree
of depth 6 needs at most $2^{6}-1 = 63$ rule choices. A 64--codon genome
therefore typically maps without wrapping, which makes the genotype to
phenotype mapping easier to reason about and limits the disruptive
effect of crossover on long genomes.

% ──────────────────────────────────────────────────────────────────────────────
\section{Fitness Function and Fitness Evaluation}

Raw fitness is the parsimony--adjusted training MSE:
\begin{equation}
f_{\text{raw}}(i) \;=\; \mathrm{MSE}(i) \,\bigl(1 + c \cdot s(i)\bigr)
\end{equation}
where $s(i)$ is the number of nodes in the phenotype AST and
$c = 10^{-3}$ a small parsimony coefficient. Parsimony pressure is
identical to that used in A1 and A2 and is included because GE, like
GP, suffers from bloat: in early experiments without parsimony the mean
phenotype size grew unboundedly after the structural sharing penalty
spread the population across many niches. Individuals whose mapping is
incomplete (more than \texttt{MAX\_WRAPS} wraps) or whose evaluation
produces NaN or infinity are assigned $f_{\text{raw}} = \infty$, which
eliminates them from selection.

The adjusted (shared) fitness used for \emph{parent selection only} is
given by Equation~\ref{eq:sharing}. The elite slot --- one slot per
generation --- is filled by the individual with the lowest raw fitness,
not the lowest adjusted fitness.

Evaluation is purely deterministic and runs the phenotype interpreter
once per training example. Protected division is used: if the
denominator has absolute value below $10^{-9}$ the operator returns 1
rather than producing a division--by--zero.

% ──────────────────────────────────────────────────────────────────────────────
\section{Selection Method}

Parent selection is tournament selection of size 2 over the adjusted
fitness. Two individuals are sampled with replacement and the one with
lower adjusted fitness is selected. The tournament size was kept at 2
for direct parity with A1 and A2; the same low selection pressure
preserves population variety long enough for the structural sharing
penalty to have a measurable effect.

% ──────────────────────────────────────────────────────────────────────────────
\section{Genetic Operators}

Three operators act on the population each generation: \emph{elitism},
\emph{one--point crossover} on integer codon arrays, and
\emph{integer mutation} per codon. Operator application uses an
\emph{exact--rate split} rather than per--child Bernoulli probabilities:
given a population of size $N$, exactly one slot is filled by elitism,
exactly $\lfloor 0.70(N-1)\rceil$ slots are produced by crossover, and the
remaining $\lfloor 0.30(N-1)\rceil$ slots are produced by mutation. The
crossover and mutation rates therefore sum to $1.0$. Each operator is
justified individually below.

\paragraph{Elitism (one slot per generation).} The single best
individual by \emph{raw} fitness is copied unchanged into the next
generation. This guarantees the monotone improvement of the
best--so--far behavioural solution from one generation to the next ---
the destabilising effect of crossover and mutation can otherwise
destroy a hard--earned good individual when the structural sharing
penalty has demoted it because its niche has become crowded. As noted
in Section~3, raw rather than adjusted fitness is used here because
elitism is about preserving the best \emph{behavioural} solution.

\paragraph{One--point crossover on codons (rate 0.70).}
A single crossover point is sampled uniformly along the 64--codon
genome; the child inherits codons before the cut from parent~1 and
codons after the cut from parent~2. This is the standard recombination
operator of canonical GE \cite{ONeillRyan2001}. It is included because
in a leftmost derivation the early codons control the structural choices
made high in the derivation tree --- the program's skeleton --- while
later codons refine sub--expressions and constants. A one--point cut
therefore cleanly exchanges the \emph{skeleton} of one parent with the
\emph{refinement} of another, which is exactly the recombination of
building blocks that drives improvement in GE. Crossover is included as
the principal \emph{exploitation} operator: it produces new individuals
entirely from material that has already survived selection, which is the
mechanism by which the search consolidates partial solutions into
stronger ones.

\paragraph{Integer mutation (per--codon probability 0.10).}
Thirty percent of children come from mutation. The mutation operator
takes one selected parent and produces a child in which each codon is
independently replaced with probability $0.10$ by a fresh uniform draw
in $[0, 256)$. Mutation is included as the principal \emph{exploration}
operator: because a single codon change can swap an entire production
choice (and therefore replace one sub--expression skeleton with another),
integrating mutation into 30\% of the new population keeps the search
from collapsing onto whatever skeleton the elite currently uses. In
pilot runs with the operator split set to $0.9/0.1$ the search consistently
converged onto the trivial predictor $\hat{y}=y_{t-1}$ in six of ten
seeds; restoring the $0.7/0.3$ balance yielded a more diverse spread of
final solutions, confirming the role of mutation as the exploratory
driver.

Together the three operators implement a balanced
exploitation/exploration cycle: elitism preserves the best behaviour,
crossover recombines existing structural building blocks, and mutation
seeds the population with previously unseen production choices.

% ──────────────────────────────────────────────────────────────────────────────
\section{Experimental Setup}

\subsection{Hardware and software environment}
All simulations were executed on the following machine:
\begin{itemize}[leftmargin=2em]
    \item Apple MacBook Pro, M--series ARM64 CPU
    \item macOS 14
    \item Oracle OpenJDK 17 (single--threaded)
    \item No external libraries; only the Java standard library is used
\end{itemize}

\subsection{Parameter values}
The GE control parameters were chosen as follows. Where a parameter
also exists in A1/A2 it was kept identical to allow direct comparison
of the search engines themselves; structural sharing parameters were
inherited from A2 because they were already validated there.

\begin{table}[H]
\centering
\begin{tabular}{l c l}
\toprule
\textbf{Parameter} & \textbf{Value} & \textbf{Justification} \\
\midrule
Population size              & 1000     & Matches A1/A2 \\
Generations                  & 100      & Matches A1/A2 \\
Tournament size              & 2        & Low pressure preserves variety \\
Crossover rate               & 0.70     & Exploitation operator (see Sec.~8) \\
Mutation rate                & 0.30     & Exploration operator; $0.70+0.30=1.0$ \\
Per--codon mutation prob.    & 0.10     & $\approx 6$ codons/genome on average \\
Genome length                & 64       & Upper bound on derivation steps at depth 6 \\
Codon range                  & $[0,256)$ & Standard 8--bit codons \\
Max derivation depth         & 6        & Matches A1/A2 GP depth cap \\
Max wraps                    & 2        & Standard GE value \\
Window size (lags)           & 7        & Matches A1/A2 \\
Train/test ratio             & 0.80     & Temporal split (no shuffling) \\
$\sigma$ (similarity thresh.) & 0.60    & Matches A2 \\
$\alpha$ (sharing coeff.)    & 0.01     & Matches A2 \\
Structural sample size       & 50       & Linear--time niche estimate \\
Parsimony coefficient        & $10^{-3}$ & Matches A1/A2 \\
Independent runs             & 10       & Seeds 101, 202, \ldots, 1010 \\
\bottomrule
\end{tabular}
\caption{SBGE control parameters used in all experiments.}
\end{table}

Two control parameters were tuned for this assignment: the operator
split and the per--codon mutation probability. The operator split was
first set to $0.9/0.1$ on the basis that crossover is GE's principal
recombination operator; however, six of ten pilot runs converged onto
the trivial $\hat{y}=y_{t-1}$ predictor under that setting because
mutation--free children dominated the new population and the elite
rapidly homogenised the population through tournament selection.
Resetting the split to $0.7/0.3$ removed this collapse mode and was
adopted. The per--codon probability was then set to $0.10$ ($\approx 6$
codon changes per mutated child) to give a meaningful structural
perturbation while keeping the average derivation--signature change
below half the genome. The remaining hyper--parameters were inherited
from A1/A2 and not re--tuned.

\subsection{Reproducibility}
The same set of ten seeds $\{101, 202, 303, 404, 505, 606, 707, 808,
909, 1010\}$ is used as in A1 and A2. With identical seeds, the same
random codon sequence is generated, so any reader who recompiles the
source code on a JDK 17 runtime will obtain the exact numerical results
reported in Section~10.

% ──────────────────────────────────────────────────────────────────────────────
\section{Results and Discussion}

\subsection{Aggregated results across 10 runs}

\begin{table}[H]
\centering
\begin{tabular}{l c c}
\toprule
\textbf{Statistic} & \textbf{Train MSE} & \textbf{Test MSE} \\
\midrule
Mean over 10 runs                & $1.9314 \times 10^{-4}$  & $2.3663 \times 10^{-4}$ \\
Standard deviation               & $2.4060 \times 10^{-5}$  & --- \\
Best run (Run 4, seed 404)       & $1.6529 \times 10^{-4}$  & $1.9213 \times 10^{-4}$ \\
\bottomrule
\end{tabular}
\caption{Summary statistics for SBGE across the ten independent runs.}
\end{table}

\subsection{Per--run results}

\begin{table}[H]
\centering
\begin{tabular}{c c c c}
\toprule
\textbf{Run} & \textbf{Train MSE} & \textbf{Test MSE} & \textbf{Time (s)} \\
\midrule
1  & 0.00018336 & 0.00021950 & 20.89 \\
2  & 0.00021701 & 0.00026954 & 23.45 \\
3  & 0.00016685 & 0.00020368 & 31.22 \\
4  & \textbf{0.00016529} & \textbf{0.00019213} & 39.77 \\
5  & 0.00021701 & 0.00026954 & 19.69 \\
6  & 0.00021607 & 0.00026763 & 21.80 \\
7  & 0.00016588 & 0.00020372 & 60.86 \\
8  & 0.00021701 & 0.00026954 & 33.53 \\
9  & 0.00016656 & 0.00020373 & 69.53 \\
10 & 0.00021639 & 0.00026731 & 45.10 \\
\midrule
Mean & 0.00019314 & 0.00023663 & 36.58 \\
\bottomrule
\end{tabular}
\caption{Per--run results across the ten seeds. Best run bolded.}
\end{table}

\subsection{Best evolved expression}
The best individual found across all ten runs was produced by Run~4
(seed 404):
\[
\hat{y}_t \;=\; \frac{2.92\,x_7}{4.48} + \frac{x_4}{1.75\,(x_7+0.10)\,x_7\,\bigl(\tfrac{8.08}{x_3} + 8.86\bigr)}.
\]
The expression combines a scaled $y_{t-7}$ baseline with a small
residual correction driven by $y_{t-3}$ and $y_{t-4}$. Although less
interpretable than a pure moving average, its test MSE of
$1.92 \times 10^{-4}$ is the closest any SBGE run came to the
Assignment~1 baseline.

\subsection{Per--generation behaviour of the best run}

\begin{table}[H]
\centering
\small
\begin{tabular}{c c c c c c}
\toprule
\textbf{Gen} & \textbf{Avg Train MSE} & \textbf{Avg Size} & \textbf{Hits} & \textbf{Variety} & \textbf{Struct.\ Sim.} \\
\midrule
10  & $1.64\!\times\!10^{3}$ & 1.75 & 39\,240 & 49  & 0.46 \\
20  & $1.22\!\times\!10^{3}$ & 2.89 & 39\,088 & 103 & 0.37 \\
30  & $1.44\!\times\!10^{2}$ & 3.78 & 39\,088 & 122 & 0.09 \\
40  & $8.45\!\times\!10^{7}$ & 4.68 & 39\,161 & 245 & 0.09 \\
50  & $5.43\!\times\!10^{4}$ & 5.61 & 39\,161 & 334 & 0.05 \\
60  & $2.09\!\times\!10^{4}$ & 5.33 & 41\,903 & 372 & 0.11 \\
70  & $5.81\!\times\!10^{2}$ & 5.39 & 42\,258 & 381 & 0.11 \\
80  & $3.22\!\times\!10^{3}$ & 5.49 & 42\,258 & 388 & 0.08 \\
90  & $3.95\!\times\!10^{6}$ & 5.23 & 42\,852 & 433 & 0.05 \\
100 & $1.18\!\times\!10^{3}$ & 5.28 & 42\,852 & 435 & 0.12 \\
\bottomrule
\end{tabular}
\caption{Per--generation statistics for the best run (Run 4, seed 404).
``Variety'' is the number of distinct phenotype expressions in the
population; ``Struct.\ Sim.'' is the mean pairwise normalised prefix
match over the derivation signatures of 50 sampled pairs.}
\end{table}

\subsection{Discussion}

\paragraph{Convergence behaviour.}
The hits count climbs from $39{,}240$ ($49\%$ of the training set) at
generation~10 to $42{,}852$ ($53.6\%$) by generation~100, while the
average phenotype size grows from $1.75$ to about $5.3$. Hits saturate
well before the MSE saturates, indicating that once the search has
found a structurally adequate skeleton further refinement happens
through small adjustments of the embedded constants rather than
through changes to the shape of the expression.

\paragraph{Why the average MSE is so large.}
The reported \emph{average} training MSE per generation is enormous
because it is dominated by a tail of individuals whose evolved
expression contains an unfortunate combination of multiplications and
mid--range constants. Such phenotypes produce outputs in the
thousands or millions, even though they are valid programs and
therefore counted in the average. The structural sharing penalty
encourages the population to maintain such structurally distinct
individuals --- which is exactly the design intent --- but it inflates
the per--generation mean. The robust measures (best--so--far raw
fitness, hits) tell the real story: the search consistently improves
its best individual generation by generation.

\paragraph{Effect of fitness sharing on diversity.}
The average pairwise structural similarity falls steadily from $0.46$
at generation~10 to between $0.05$ and $0.12$ from generation~50
onwards, and phenotype variety rises from $49$ to over $400$.
Compared with the equivalent figures recorded for the canonical
Assignment~1 GP run (variety collapsed to the low hundreds by
generation~100), the sharing penalty is clearly maintaining a wider
structural spread throughout the search. This is the expected
qualitative behaviour of fitness--sharing in any evolutionary
algorithm, and it is reproduced here on derivation signatures rather
than on tree topology.

\paragraph{Success predicate.}
Five of the ten runs reach the success threshold of
$\mathrm{MSE} < 2\times 10^{-4}$. The remaining five runs (2, 5, 6, 8,
10) all converge to the same plateau ($\mathrm{MSE}\approx
2.17\times 10^{-4}$), which corresponds to the trivial single--lag
predictor $\hat{y}=y_{t-1}$. Inspection of those runs' per--generation
statistics shows that they exhausted the early structural diversity
before the sharing penalty could move them off this attractor: in GE,
as in canonical GP, the most recent lag is a deep local optimum, and a
stronger exploration mechanism (e.g.\ random restarts or grammar
immunisation) would likely be required to consistently escape it.

% ──────────────────────────────────────────────────────────────────────────────
\section{Runtimes}

The average wall--clock time per run across the ten seeds was
$36.58\,\text{s}$, with a minimum of $19.69\,\text{s}$ (Run~5) and a
maximum of $69.53\,\text{s}$ (Run~9). The total wall--clock time for
all ten runs was approximately $6\,\text{min}\,06\,\text{s}$. The
variance in per--run time reflects the variance in average phenotype
size: runs that grow longer expressions spend proportionally more time
in the per--example interpreter.

GE is substantially faster than the Assignment~2 SBGP engine (mean
$671\,\text{s}$ per run) for two reasons. First, the GE phenotype is
already represented as a compact binary AST built directly during
mapping, whereas A2's GP engine built a richer \texttt{Node} tree
carrying structural metadata. Second, derivation--signature similarity
is computed on a flat \texttt{int[]} via a single tight loop, while
A2's prefix matching on subtree nodes followed an object--graph
recursion. Both effects compound when the structural sample size is
$50$ per individual per generation.

% ──────────────────────────────────────────────────────────────────────────────
\section{Comparison with Assignments 1 and 2}

\begin{table}[H]
\centering
\resizebox{\textwidth}{!}{%
\begin{tabular}{l c c c c}
\toprule
\textbf{Method} & \textbf{Avg Train MSE} & \textbf{Std.\ Dev.} & \textbf{Best Train MSE} & \textbf{Avg Test MSE} \\
\midrule
A1 (Canonical GP)             & $1.4812\!\times\!10^{-4}$ & $0.66\!\times\!10^{-5}$ & $1.4257\!\times\!10^{-4}$ & $1.7889\!\times\!10^{-4}$ \\
A2 (SBGP, fitness sharing)    & $1.4611\!\times\!10^{-4}$ & $0.34\!\times\!10^{-5}$ & $1.4120\!\times\!10^{-4}$ & $1.7394\!\times\!10^{-4}$ \\
A3 (SBGE, this work)          & $1.9314\!\times\!10^{-4}$ & $2.41\!\times\!10^{-5}$ & $1.6529\!\times\!10^{-4}$ & $2.3663\!\times\!10^{-4}$ \\
\bottomrule
\end{tabular}}
\caption{End--to--end comparison across the three assignments. All three
were trained on the same $79{,}994$ examples and tested on the same
$19{,}999$ examples with identical seeds.}
\end{table}

\subsection{Quality of solutions}
A1 and A2 (both tree--GP) produce noticeably lower training and test
MSE than A3 (GE). The A1 canonical GP attains an average training MSE
$\sim23\%$ lower than A3 and the A2 SBGP attains $\sim24\%$ lower. The
test--set ranking is the same: A2 $<$ A1 $<$ A3. The best individual
produced by A3 (Run 4, train MSE $1.65\times 10^{-4}$, test MSE
$1.92\times 10^{-4}$) does, however, come within $16\%$ of the A1 best
and outperforms the A1 \emph{average}, showing that GE is capable of
finding competitive solutions when its search trajectory survives the
single--lag attractor described in Section~10.5. This pattern is
consistent with the wider literature: on small symbolic regression
problems canonical GE typically underperforms tree--GP because (a)
grammar mapping introduces a layer of indirection between genotype and
phenotype, so a single codon change can disrupt an arbitrary subtree
rather than a single subtree node; and (b) one--point crossover on
codons is structurally less informed than subtree crossover.

\subsection{Robustness}
The standard deviation of training MSE is $\approx 3.7\times$ larger
for A3 than for A1 and $\approx 7\times$ larger than for A2. This
greater run--to--run variance is again attributable to the
codon--to--tree mapping: a small early bias in which production codons
appear in the initial population can lock the search onto a particular
sub--region of program space, and the linear genome operators are not
as effective as subtree crossover at escaping that bias. A2's
structural sharing penalty was specifically able to reduce the A1
variance by half because it operated directly on the GP tree; the same
penalty in A3 operates one indirection step away and is therefore less
effective in absolute terms.

\subsection{Speed}
A3 is by far the fastest of the three: average $35\,\text{s}$ per run
versus $298\,\text{s}$ for A1 and $671\,\text{s}$ for A2. Per the
explanation in Section~11, GE pays for this speed with a compact AST
representation and an efficient signature--matching routine.

\subsection{Summary}
The three assignments together produce a clean experimental story.
Tree--GP (A1) provides a strong baseline. Adding structural fitness
sharing to GP (A2) gives a small but consistent improvement in both
mean quality and run--to--run robustness. Replacing the search engine
with canonical GE and applying the same structural sharing idea (A3)
trades some predictive accuracy for substantially lower wall--clock
time and a more compact, interpretable best solution. For the
electricity--load regression task the GP family remains the better
choice when accuracy is the priority; SBGE is the better choice when
training time is the priority or when a grammar--constrained search
space is required for downstream interpretability.

% ──────────────────────────────────────────────────────────────────────────────
\begin{thebibliography}{9}
\bibitem{ONeillRyan2001}
M.~O'Neill and C.~Ryan.
\newblock Grammatical Evolution.
\newblock \emph{IEEE Transactions on Evolutionary Computation}, 5(4):349--358, 2001.

\bibitem{RyanAzad2003}
R.~M.~A.~Azad and C.~Ryan.
\newblock Sensible Initialisation in Chorus.
\newblock In \emph{Proc.\ EuroGP}, pages 394--403, Springer, 2003.

\bibitem{PonyGE2}
M.~Fenton, J.~McDermott, D.~Fagan, S.~Forstenlechner, M.~O'Neill, and
E.~Hemberg.
\newblock PonyGE2: Grammatical Evolution in Python.
\newblock In \emph{GECCO Companion}, pages 1194--1201, ACM, 2017.
\end{thebibliography}

\end{document}
